% Felles oppsett for kretsskjemaene i Fysikk 2. Brukes med % Felles oppsett for kretsskjemaene i Fysikk 2. Brukes med % Felles oppsett for kretsskjemaene i Fysikk 2. Brukes med % Felles oppsett for kretsskjemaene i Fysikk 2. Brukes med \input{krets}.
% Bygger på felles.tex (farger: acc -> --accent, mut -> --muted, hi -> --st-forstatt).
% Hvitt fyll (åpne terminaler, bryterkontakter) blir --bg i build.py, så det følger temaet.
\input{felles}
\usepackage[american]{circuitikz}
\ctikzset{
  bipoles/thickness=1.2,
  bipoles/length=1.15cm,
  resistors/scale=0.85,
  current arrow scale=11,
  voltage/distance from node=.45,
  voltage/bump b=1.6,
}
\tikzset{
  every picture/.append style={line width=0.6pt, font=\small},
  % maskestrøm: sirkulær pil med klokka rundt (x,y)
  maske/.style={acc, -{Stealth[length=4pt]}, line width=0.7pt},
  node/.style={circle, fill=acc, inner sep=1.3pt},
  nodelbl/.style={text=acc, font=\small},
}
% \maske{x}{y}{navn}: maskestrøm med klokka, sentrert i (x,y)
\newcommand{\maske}[3]{%
  \draw[maske] ($(#1,#2)+(60:0.42)$) arc[start angle=60, end angle=-240, radius=0.42];
  \node[text=acc, font=\small] at (#1,#2) {#3};}

% Konvensjoner (testet):
%   V: + ved startpunktet, så en kilde tegnet oppover trenger invert for å få + oppe.
%   I: pila peker i tegneretningen.
%   Loddrett oppover: l = venstre, a = høyre.  Loddrett ned: l_ = venstre, a^ = høyre.
%   Vannrett: l = over, a = under.
% \volt{x}{ytopp}{ybunn}{etikett}: + oppe, - nede, etikett midt mellom, i x
\newcommand{\volt}[4]{%
  \node[font=\small] at (#1,#2) {$+$}; \node[font=\small] at (#1,#3) {$-$};
  \node[font=\small] at (#1,{(#2+#3)/2}) {#4};}
% \volth{y}{xvenstre}{xhoyre}{etikett}: + til venstre, - til høyre
\newcommand{\volth}[4]{%
  \node[font=\small] at (#2,#1) {$+$}; \node[font=\small] at (#3,#1) {$-$};
  \node[font=\small] at ({(#2+#3)/2},#1) {#4};}
\newcommand{\kO}{\,\mathrm{k}\Omega}
\newcommand{\Om}{\,\Omega}

% \thev{ETh}{RTh}: Thévenin-ekvivalent med terminaler a og b, origo nede til venstre
\newcommand{\thev}[2]{%
  \draw (0,0) to[V, invert, l=$E_{Th}$, a=#1] (0,2.6) to[R, l=$R_{Th}$, a=#2] (2.8,2.6)
    to[short, -o] (3.6,2.6) node[above] {$a$};
  \draw (0,0) to[short, -o] (3.6,0) node[below] {$b$};}
% \nort{IN}{RN}: Norton-ekvivalent med terminaler a og b
\newcommand{\nort}[2]{%
  \draw (0,0) to[I, l=$I_N$, a=#1] (0,2.6) -- (2.6,2.6) to[short, -o] (3.6,2.6) node[above] {$a$};
  \draw (2.6,2.6) to[R, l_=$R_N$, a^=#2, *-*] (2.6,0);
  \draw (0,0) to[short, -o] (3.6,0) node[below] {$b$};}
% \last{navn}{verdi}: lastmotstand mellom terminalene a og b (x = 3.6), i aksentfarge
\newcommand{\last}[2]{%
  \draw (3.6,2.6) -- (4.8,2.6) to[R, l_=#1, a^=#2, color=acc] (4.8,0) -- (3.6,0);}

% Transientgrafer (pgfplots): x i tidskonstanter eller sekunder, ingen rutenett
\pgfplotsset{
  trans/.style={width=7.4cm, height=4.6cm, axis lines=left, clip=false, samples=80,
    tick label style={font=\small}, label style={font=\small}, scaled ticks=false,
    /pgf/number format/1000 sep={}, every axis plot/.append style={line width=0.8pt},
    xlabel style={at={(axis description cs:1.03,0)}, anchor=west},
    ylabel style={rotate=-90, at={(axis description cs:0,1.04)}, anchor=south}},
}

% Strømbaner for animasjon (strom.js): f01 ... f12 tegnes i strømretningen og blir
% CSS-klassene flow f01 ... i build.py. Tegnes sist, slik at prikkene ligger over ledningene.
\definecolor{f01}{HTML}{F0F001}
\definecolor{f02}{HTML}{F0F002}
\definecolor{f03}{HTML}{F0F003}
\definecolor{f04}{HTML}{F0F004}
\definecolor{f05}{HTML}{F0F005}
\definecolor{f06}{HTML}{F0F006}
\definecolor{f07}{HTML}{F0F007}
\definecolor{f08}{HTML}{F0F008}
\definecolor{f09}{HTML}{F0F009}
\definecolor{f10}{HTML}{F0F010}
\definecolor{f11}{HTML}{F0F011}
\definecolor{f12}{HTML}{F0F012}
\tikzset{strom/.style={line width=1.6pt, line cap=round, line join=round}}
.
% Bygger på felles.tex (farger: acc -> --accent, mut -> --muted, hi -> --st-forstatt).
% Hvitt fyll (åpne terminaler, bryterkontakter) blir --bg i build.py, så det følger temaet.
% Felles oppsett for figurene. Brukes med \input{felles} etter \documentclass.
% Fargene er plassholdere som build.py bytter til CSS-variabler:
%   svart -> tekstfarge, acc -> --accent, mut -> --muted, hi -> --st-forstatt
\usepackage{amsmath}
\usepackage{pgfplots}
\pgfplotsset{compat=1.18}
\usetikzlibrary{arrows.meta, positioning, calc}
\definecolor{acc}{HTML}{FF0000}
\definecolor{mut}{HTML}{00FF00}
\definecolor{hi}{HTML}{0000FF}

% Normaltetthet med forventning #1 og standardavvik #2
\pgfmathdeclarefunction{gauss}{2}{%
  \pgfmathparse{exp(-((x-#1)^2)/(2*#2^2))/(#2*sqrt(2*pi))}}
% t-tetthet med #1 frihetsgrader; #2 er normeringskonstanten (regnet ut i normal.py)
\pgfmathdeclarefunction{tdens}{2}{%
  \pgfmathparse{#2*(1+x^2/#1)^(-(#1+1)/2)}}
% Eksponentialtetthet med rate #1
\pgfmathdeclarefunction{expdens}{1}{\pgfmathparse{#1*exp(-#1*x)}}

\tikzset{
  % hendelsestre som vokser mot høyre
  tre/.style={grow'=right, edge from parent/.style={draw, line width=0.5pt}},
  p/.style={font=\small, text=acc, inner sep=2pt},
  up/.style={p, above, sloped}, dn/.style={p, below, sloped},
  leaf/.style={anchor=west, xshift=0.35cm},
  rot/.style={circle, fill, inner sep=1.4pt},
  hjelp/.style={mut, dashed, line width=0.5pt},
}

\pgfplotsset{
  % kurve uten y-akse: bare x-aksen med de merkene figuren trenger
  kurve/.style={
    width=11cm, height=4.6cm, axis x line*=bottom, hide y axis,
    ymin=0, enlarge x limits=false, clip=false, samples=120,
    xtick style={draw=none}, tick label style={font=\small}, scaled x ticks=false,
    every axis plot/.append style={line width=0.7pt}},
  % spredningsplott med vanlige akser
  plott/.style={
    width=9cm, height=5.6cm, axis lines=left, clip mode=individual,
    tick label style={font=\small}, label style={font=\small},
    scaled ticks=false, /pgf/number format/1000 sep={\,},
    enlarge x limits=0.02, enlarge y limits=0.08},
  areal/.style={draw=none, fill=acc, fill opacity=0.3},
  areal2/.style={draw=none, fill=hi, fill opacity=0.3},
}

\usepackage[american]{circuitikz}
\ctikzset{
  bipoles/thickness=1.2,
  bipoles/length=1.15cm,
  resistors/scale=0.85,
  current arrow scale=11,
  voltage/distance from node=.45,
  voltage/bump b=1.6,
}
\tikzset{
  every picture/.append style={line width=0.6pt, font=\small},
  % maskestrøm: sirkulær pil med klokka rundt (x,y)
  maske/.style={acc, -{Stealth[length=4pt]}, line width=0.7pt},
  node/.style={circle, fill=acc, inner sep=1.3pt},
  nodelbl/.style={text=acc, font=\small},
}
% \maske{x}{y}{navn}: maskestrøm med klokka, sentrert i (x,y)
\newcommand{\maske}[3]{%
  \draw[maske] ($(#1,#2)+(60:0.42)$) arc[start angle=60, end angle=-240, radius=0.42];
  \node[text=acc, font=\small] at (#1,#2) {#3};}

% Konvensjoner (testet):
%   V: + ved startpunktet, så en kilde tegnet oppover trenger invert for å få + oppe.
%   I: pila peker i tegneretningen.
%   Loddrett oppover: l = venstre, a = høyre.  Loddrett ned: l_ = venstre, a^ = høyre.
%   Vannrett: l = over, a = under.
% \volt{x}{ytopp}{ybunn}{etikett}: + oppe, - nede, etikett midt mellom, i x
\newcommand{\volt}[4]{%
  \node[font=\small] at (#1,#2) {$+$}; \node[font=\small] at (#1,#3) {$-$};
  \node[font=\small] at (#1,{(#2+#3)/2}) {#4};}
% \volth{y}{xvenstre}{xhoyre}{etikett}: + til venstre, - til høyre
\newcommand{\volth}[4]{%
  \node[font=\small] at (#2,#1) {$+$}; \node[font=\small] at (#3,#1) {$-$};
  \node[font=\small] at ({(#2+#3)/2},#1) {#4};}
\newcommand{\kO}{\,\mathrm{k}\Omega}
\newcommand{\Om}{\,\Omega}

% \thev{ETh}{RTh}: Thévenin-ekvivalent med terminaler a og b, origo nede til venstre
\newcommand{\thev}[2]{%
  \draw (0,0) to[V, invert, l=$E_{Th}$, a=#1] (0,2.6) to[R, l=$R_{Th}$, a=#2] (2.8,2.6)
    to[short, -o] (3.6,2.6) node[above] {$a$};
  \draw (0,0) to[short, -o] (3.6,0) node[below] {$b$};}
% \nort{IN}{RN}: Norton-ekvivalent med terminaler a og b
\newcommand{\nort}[2]{%
  \draw (0,0) to[I, l=$I_N$, a=#1] (0,2.6) -- (2.6,2.6) to[short, -o] (3.6,2.6) node[above] {$a$};
  \draw (2.6,2.6) to[R, l_=$R_N$, a^=#2, *-*] (2.6,0);
  \draw (0,0) to[short, -o] (3.6,0) node[below] {$b$};}
% \last{navn}{verdi}: lastmotstand mellom terminalene a og b (x = 3.6), i aksentfarge
\newcommand{\last}[2]{%
  \draw (3.6,2.6) -- (4.8,2.6) to[R, l_=#1, a^=#2, color=acc] (4.8,0) -- (3.6,0);}

% Transientgrafer (pgfplots): x i tidskonstanter eller sekunder, ingen rutenett
\pgfplotsset{
  trans/.style={width=7.4cm, height=4.6cm, axis lines=left, clip=false, samples=80,
    tick label style={font=\small}, label style={font=\small}, scaled ticks=false,
    /pgf/number format/1000 sep={}, every axis plot/.append style={line width=0.8pt},
    xlabel style={at={(axis description cs:1.03,0)}, anchor=west},
    ylabel style={rotate=-90, at={(axis description cs:0,1.04)}, anchor=south}},
}

% Strømbaner for animasjon (strom.js): f01 ... f12 tegnes i strømretningen og blir
% CSS-klassene flow f01 ... i build.py. Tegnes sist, slik at prikkene ligger over ledningene.
\definecolor{f01}{HTML}{F0F001}
\definecolor{f02}{HTML}{F0F002}
\definecolor{f03}{HTML}{F0F003}
\definecolor{f04}{HTML}{F0F004}
\definecolor{f05}{HTML}{F0F005}
\definecolor{f06}{HTML}{F0F006}
\definecolor{f07}{HTML}{F0F007}
\definecolor{f08}{HTML}{F0F008}
\definecolor{f09}{HTML}{F0F009}
\definecolor{f10}{HTML}{F0F010}
\definecolor{f11}{HTML}{F0F011}
\definecolor{f12}{HTML}{F0F012}
\tikzset{strom/.style={line width=1.6pt, line cap=round, line join=round}}
.
% Bygger på felles.tex (farger: acc -> --accent, mut -> --muted, hi -> --st-forstatt).
% Hvitt fyll (åpne terminaler, bryterkontakter) blir --bg i build.py, så det følger temaet.
% Felles oppsett for figurene. Brukes med % Felles oppsett for figurene. Brukes med \input{felles} etter \documentclass.
% Fargene er plassholdere som build.py bytter til CSS-variabler:
%   svart -> tekstfarge, acc -> --accent, mut -> --muted, hi -> --st-forstatt
\usepackage{amsmath}
\usepackage{pgfplots}
\pgfplotsset{compat=1.18}
\usetikzlibrary{arrows.meta, positioning, calc}
\definecolor{acc}{HTML}{FF0000}
\definecolor{mut}{HTML}{00FF00}
\definecolor{hi}{HTML}{0000FF}

% Normaltetthet med forventning #1 og standardavvik #2
\pgfmathdeclarefunction{gauss}{2}{%
  \pgfmathparse{exp(-((x-#1)^2)/(2*#2^2))/(#2*sqrt(2*pi))}}
% t-tetthet med #1 frihetsgrader; #2 er normeringskonstanten (regnet ut i normal.py)
\pgfmathdeclarefunction{tdens}{2}{%
  \pgfmathparse{#2*(1+x^2/#1)^(-(#1+1)/2)}}
% Eksponentialtetthet med rate #1
\pgfmathdeclarefunction{expdens}{1}{\pgfmathparse{#1*exp(-#1*x)}}

\tikzset{
  % hendelsestre som vokser mot høyre
  tre/.style={grow'=right, edge from parent/.style={draw, line width=0.5pt}},
  p/.style={font=\small, text=acc, inner sep=2pt},
  up/.style={p, above, sloped}, dn/.style={p, below, sloped},
  leaf/.style={anchor=west, xshift=0.35cm},
  rot/.style={circle, fill, inner sep=1.4pt},
  hjelp/.style={mut, dashed, line width=0.5pt},
}

\pgfplotsset{
  % kurve uten y-akse: bare x-aksen med de merkene figuren trenger
  kurve/.style={
    width=11cm, height=4.6cm, axis x line*=bottom, hide y axis,
    ymin=0, enlarge x limits=false, clip=false, samples=120,
    xtick style={draw=none}, tick label style={font=\small}, scaled x ticks=false,
    every axis plot/.append style={line width=0.7pt}},
  % spredningsplott med vanlige akser
  plott/.style={
    width=9cm, height=5.6cm, axis lines=left, clip mode=individual,
    tick label style={font=\small}, label style={font=\small},
    scaled ticks=false, /pgf/number format/1000 sep={\,},
    enlarge x limits=0.02, enlarge y limits=0.08},
  areal/.style={draw=none, fill=acc, fill opacity=0.3},
  areal2/.style={draw=none, fill=hi, fill opacity=0.3},
}
 etter \documentclass.
% Fargene er plassholdere som build.py bytter til CSS-variabler:
%   svart -> tekstfarge, acc -> --accent, mut -> --muted, hi -> --st-forstatt
\usepackage{amsmath}
\usepackage{pgfplots}
\pgfplotsset{compat=1.18}
\usetikzlibrary{arrows.meta, positioning, calc}
\definecolor{acc}{HTML}{FF0000}
\definecolor{mut}{HTML}{00FF00}
\definecolor{hi}{HTML}{0000FF}

% Normaltetthet med forventning #1 og standardavvik #2
\pgfmathdeclarefunction{gauss}{2}{%
  \pgfmathparse{exp(-((x-#1)^2)/(2*#2^2))/(#2*sqrt(2*pi))}}
% t-tetthet med #1 frihetsgrader; #2 er normeringskonstanten (regnet ut i normal.py)
\pgfmathdeclarefunction{tdens}{2}{%
  \pgfmathparse{#2*(1+x^2/#1)^(-(#1+1)/2)}}
% Eksponentialtetthet med rate #1
\pgfmathdeclarefunction{expdens}{1}{\pgfmathparse{#1*exp(-#1*x)}}

\tikzset{
  % hendelsestre som vokser mot høyre
  tre/.style={grow'=right, edge from parent/.style={draw, line width=0.5pt}},
  p/.style={font=\small, text=acc, inner sep=2pt},
  up/.style={p, above, sloped}, dn/.style={p, below, sloped},
  leaf/.style={anchor=west, xshift=0.35cm},
  rot/.style={circle, fill, inner sep=1.4pt},
  hjelp/.style={mut, dashed, line width=0.5pt},
}

\pgfplotsset{
  % kurve uten y-akse: bare x-aksen med de merkene figuren trenger
  kurve/.style={
    width=11cm, height=4.6cm, axis x line*=bottom, hide y axis,
    ymin=0, enlarge x limits=false, clip=false, samples=120,
    xtick style={draw=none}, tick label style={font=\small}, scaled x ticks=false,
    every axis plot/.append style={line width=0.7pt}},
  % spredningsplott med vanlige akser
  plott/.style={
    width=9cm, height=5.6cm, axis lines=left, clip mode=individual,
    tick label style={font=\small}, label style={font=\small},
    scaled ticks=false, /pgf/number format/1000 sep={\,},
    enlarge x limits=0.02, enlarge y limits=0.08},
  areal/.style={draw=none, fill=acc, fill opacity=0.3},
  areal2/.style={draw=none, fill=hi, fill opacity=0.3},
}

\usepackage[american]{circuitikz}
\ctikzset{
  bipoles/thickness=1.2,
  bipoles/length=1.15cm,
  resistors/scale=0.85,
  current arrow scale=11,
  voltage/distance from node=.45,
  voltage/bump b=1.6,
}
\tikzset{
  every picture/.append style={line width=0.6pt, font=\small},
  % maskestrøm: sirkulær pil med klokka rundt (x,y)
  maske/.style={acc, -{Stealth[length=4pt]}, line width=0.7pt},
  node/.style={circle, fill=acc, inner sep=1.3pt},
  nodelbl/.style={text=acc, font=\small},
}
% \maske{x}{y}{navn}: maskestrøm med klokka, sentrert i (x,y)
\newcommand{\maske}[3]{%
  \draw[maske] ($(#1,#2)+(60:0.42)$) arc[start angle=60, end angle=-240, radius=0.42];
  \node[text=acc, font=\small] at (#1,#2) {#3};}

% Konvensjoner (testet):
%   V: + ved startpunktet, så en kilde tegnet oppover trenger invert for å få + oppe.
%   I: pila peker i tegneretningen.
%   Loddrett oppover: l = venstre, a = høyre.  Loddrett ned: l_ = venstre, a^ = høyre.
%   Vannrett: l = over, a = under.
% \volt{x}{ytopp}{ybunn}{etikett}: + oppe, - nede, etikett midt mellom, i x
\newcommand{\volt}[4]{%
  \node[font=\small] at (#1,#2) {$+$}; \node[font=\small] at (#1,#3) {$-$};
  \node[font=\small] at (#1,{(#2+#3)/2}) {#4};}
% \volth{y}{xvenstre}{xhoyre}{etikett}: + til venstre, - til høyre
\newcommand{\volth}[4]{%
  \node[font=\small] at (#2,#1) {$+$}; \node[font=\small] at (#3,#1) {$-$};
  \node[font=\small] at ({(#2+#3)/2},#1) {#4};}
\newcommand{\kO}{\,\mathrm{k}\Omega}
\newcommand{\Om}{\,\Omega}

% \thev{ETh}{RTh}: Thévenin-ekvivalent med terminaler a og b, origo nede til venstre
\newcommand{\thev}[2]{%
  \draw (0,0) to[V, invert, l=$E_{Th}$, a=#1] (0,2.6) to[R, l=$R_{Th}$, a=#2] (2.8,2.6)
    to[short, -o] (3.6,2.6) node[above] {$a$};
  \draw (0,0) to[short, -o] (3.6,0) node[below] {$b$};}
% \nort{IN}{RN}: Norton-ekvivalent med terminaler a og b
\newcommand{\nort}[2]{%
  \draw (0,0) to[I, l=$I_N$, a=#1] (0,2.6) -- (2.6,2.6) to[short, -o] (3.6,2.6) node[above] {$a$};
  \draw (2.6,2.6) to[R, l_=$R_N$, a^=#2, *-*] (2.6,0);
  \draw (0,0) to[short, -o] (3.6,0) node[below] {$b$};}
% \last{navn}{verdi}: lastmotstand mellom terminalene a og b (x = 3.6), i aksentfarge
\newcommand{\last}[2]{%
  \draw (3.6,2.6) -- (4.8,2.6) to[R, l_=#1, a^=#2, color=acc] (4.8,0) -- (3.6,0);}

% Transientgrafer (pgfplots): x i tidskonstanter eller sekunder, ingen rutenett
\pgfplotsset{
  trans/.style={width=7.4cm, height=4.6cm, axis lines=left, clip=false, samples=80,
    tick label style={font=\small}, label style={font=\small}, scaled ticks=false,
    /pgf/number format/1000 sep={}, every axis plot/.append style={line width=0.8pt},
    xlabel style={at={(axis description cs:1.03,0)}, anchor=west},
    ylabel style={rotate=-90, at={(axis description cs:0,1.04)}, anchor=south}},
}

% Strømbaner for animasjon (strom.js): f01 ... f12 tegnes i strømretningen og blir
% CSS-klassene flow f01 ... i build.py. Tegnes sist, slik at prikkene ligger over ledningene.
\definecolor{f01}{HTML}{F0F001}
\definecolor{f02}{HTML}{F0F002}
\definecolor{f03}{HTML}{F0F003}
\definecolor{f04}{HTML}{F0F004}
\definecolor{f05}{HTML}{F0F005}
\definecolor{f06}{HTML}{F0F006}
\definecolor{f07}{HTML}{F0F007}
\definecolor{f08}{HTML}{F0F008}
\definecolor{f09}{HTML}{F0F009}
\definecolor{f10}{HTML}{F0F010}
\definecolor{f11}{HTML}{F0F011}
\definecolor{f12}{HTML}{F0F012}
\tikzset{strom/.style={line width=1.6pt, line cap=round, line join=round}}
.
% Bygger på felles.tex (farger: acc -> --accent, mut -> --muted, hi -> --st-forstatt).
% Hvitt fyll (åpne terminaler, bryterkontakter) blir --bg i build.py, så det følger temaet.
% Felles oppsett for figurene. Brukes med % Felles oppsett for figurene. Brukes med % Felles oppsett for figurene. Brukes med \input{felles} etter \documentclass.
% Fargene er plassholdere som build.py bytter til CSS-variabler:
%   svart -> tekstfarge, acc -> --accent, mut -> --muted, hi -> --st-forstatt
\usepackage{amsmath}
\usepackage{pgfplots}
\pgfplotsset{compat=1.18}
\usetikzlibrary{arrows.meta, positioning, calc}
\definecolor{acc}{HTML}{FF0000}
\definecolor{mut}{HTML}{00FF00}
\definecolor{hi}{HTML}{0000FF}

% Normaltetthet med forventning #1 og standardavvik #2
\pgfmathdeclarefunction{gauss}{2}{%
  \pgfmathparse{exp(-((x-#1)^2)/(2*#2^2))/(#2*sqrt(2*pi))}}
% t-tetthet med #1 frihetsgrader; #2 er normeringskonstanten (regnet ut i normal.py)
\pgfmathdeclarefunction{tdens}{2}{%
  \pgfmathparse{#2*(1+x^2/#1)^(-(#1+1)/2)}}
% Eksponentialtetthet med rate #1
\pgfmathdeclarefunction{expdens}{1}{\pgfmathparse{#1*exp(-#1*x)}}

\tikzset{
  % hendelsestre som vokser mot høyre
  tre/.style={grow'=right, edge from parent/.style={draw, line width=0.5pt}},
  p/.style={font=\small, text=acc, inner sep=2pt},
  up/.style={p, above, sloped}, dn/.style={p, below, sloped},
  leaf/.style={anchor=west, xshift=0.35cm},
  rot/.style={circle, fill, inner sep=1.4pt},
  hjelp/.style={mut, dashed, line width=0.5pt},
}

\pgfplotsset{
  % kurve uten y-akse: bare x-aksen med de merkene figuren trenger
  kurve/.style={
    width=11cm, height=4.6cm, axis x line*=bottom, hide y axis,
    ymin=0, enlarge x limits=false, clip=false, samples=120,
    xtick style={draw=none}, tick label style={font=\small}, scaled x ticks=false,
    every axis plot/.append style={line width=0.7pt}},
  % spredningsplott med vanlige akser
  plott/.style={
    width=9cm, height=5.6cm, axis lines=left, clip mode=individual,
    tick label style={font=\small}, label style={font=\small},
    scaled ticks=false, /pgf/number format/1000 sep={\,},
    enlarge x limits=0.02, enlarge y limits=0.08},
  areal/.style={draw=none, fill=acc, fill opacity=0.3},
  areal2/.style={draw=none, fill=hi, fill opacity=0.3},
}
 etter \documentclass.
% Fargene er plassholdere som build.py bytter til CSS-variabler:
%   svart -> tekstfarge, acc -> --accent, mut -> --muted, hi -> --st-forstatt
\usepackage{amsmath}
\usepackage{pgfplots}
\pgfplotsset{compat=1.18}
\usetikzlibrary{arrows.meta, positioning, calc}
\definecolor{acc}{HTML}{FF0000}
\definecolor{mut}{HTML}{00FF00}
\definecolor{hi}{HTML}{0000FF}

% Normaltetthet med forventning #1 og standardavvik #2
\pgfmathdeclarefunction{gauss}{2}{%
  \pgfmathparse{exp(-((x-#1)^2)/(2*#2^2))/(#2*sqrt(2*pi))}}
% t-tetthet med #1 frihetsgrader; #2 er normeringskonstanten (regnet ut i normal.py)
\pgfmathdeclarefunction{tdens}{2}{%
  \pgfmathparse{#2*(1+x^2/#1)^(-(#1+1)/2)}}
% Eksponentialtetthet med rate #1
\pgfmathdeclarefunction{expdens}{1}{\pgfmathparse{#1*exp(-#1*x)}}

\tikzset{
  % hendelsestre som vokser mot høyre
  tre/.style={grow'=right, edge from parent/.style={draw, line width=0.5pt}},
  p/.style={font=\small, text=acc, inner sep=2pt},
  up/.style={p, above, sloped}, dn/.style={p, below, sloped},
  leaf/.style={anchor=west, xshift=0.35cm},
  rot/.style={circle, fill, inner sep=1.4pt},
  hjelp/.style={mut, dashed, line width=0.5pt},
}

\pgfplotsset{
  % kurve uten y-akse: bare x-aksen med de merkene figuren trenger
  kurve/.style={
    width=11cm, height=4.6cm, axis x line*=bottom, hide y axis,
    ymin=0, enlarge x limits=false, clip=false, samples=120,
    xtick style={draw=none}, tick label style={font=\small}, scaled x ticks=false,
    every axis plot/.append style={line width=0.7pt}},
  % spredningsplott med vanlige akser
  plott/.style={
    width=9cm, height=5.6cm, axis lines=left, clip mode=individual,
    tick label style={font=\small}, label style={font=\small},
    scaled ticks=false, /pgf/number format/1000 sep={\,},
    enlarge x limits=0.02, enlarge y limits=0.08},
  areal/.style={draw=none, fill=acc, fill opacity=0.3},
  areal2/.style={draw=none, fill=hi, fill opacity=0.3},
}
 etter \documentclass.
% Fargene er plassholdere som build.py bytter til CSS-variabler:
%   svart -> tekstfarge, acc -> --accent, mut -> --muted, hi -> --st-forstatt
\usepackage{amsmath}
\usepackage{pgfplots}
\pgfplotsset{compat=1.18}
\usetikzlibrary{arrows.meta, positioning, calc}
\definecolor{acc}{HTML}{FF0000}
\definecolor{mut}{HTML}{00FF00}
\definecolor{hi}{HTML}{0000FF}

% Normaltetthet med forventning #1 og standardavvik #2
\pgfmathdeclarefunction{gauss}{2}{%
  \pgfmathparse{exp(-((x-#1)^2)/(2*#2^2))/(#2*sqrt(2*pi))}}
% t-tetthet med #1 frihetsgrader; #2 er normeringskonstanten (regnet ut i normal.py)
\pgfmathdeclarefunction{tdens}{2}{%
  \pgfmathparse{#2*(1+x^2/#1)^(-(#1+1)/2)}}
% Eksponentialtetthet med rate #1
\pgfmathdeclarefunction{expdens}{1}{\pgfmathparse{#1*exp(-#1*x)}}

\tikzset{
  % hendelsestre som vokser mot høyre
  tre/.style={grow'=right, edge from parent/.style={draw, line width=0.5pt}},
  p/.style={font=\small, text=acc, inner sep=2pt},
  up/.style={p, above, sloped}, dn/.style={p, below, sloped},
  leaf/.style={anchor=west, xshift=0.35cm},
  rot/.style={circle, fill, inner sep=1.4pt},
  hjelp/.style={mut, dashed, line width=0.5pt},
}

\pgfplotsset{
  % kurve uten y-akse: bare x-aksen med de merkene figuren trenger
  kurve/.style={
    width=11cm, height=4.6cm, axis x line*=bottom, hide y axis,
    ymin=0, enlarge x limits=false, clip=false, samples=120,
    xtick style={draw=none}, tick label style={font=\small}, scaled x ticks=false,
    every axis plot/.append style={line width=0.7pt}},
  % spredningsplott med vanlige akser
  plott/.style={
    width=9cm, height=5.6cm, axis lines=left, clip mode=individual,
    tick label style={font=\small}, label style={font=\small},
    scaled ticks=false, /pgf/number format/1000 sep={\,},
    enlarge x limits=0.02, enlarge y limits=0.08},
  areal/.style={draw=none, fill=acc, fill opacity=0.3},
  areal2/.style={draw=none, fill=hi, fill opacity=0.3},
}

\usepackage[american]{circuitikz}
\ctikzset{
  bipoles/thickness=1.2,
  bipoles/length=1.15cm,
  resistors/scale=0.85,
  current arrow scale=11,
  voltage/distance from node=.45,
  voltage/bump b=1.6,
}
\tikzset{
  every picture/.append style={line width=0.6pt, font=\small},
  % maskestrøm: sirkulær pil med klokka rundt (x,y)
  maske/.style={acc, -{Stealth[length=4pt]}, line width=0.7pt},
  node/.style={circle, fill=acc, inner sep=1.3pt},
  nodelbl/.style={text=acc, font=\small},
}
% \maske{x}{y}{navn}: maskestrøm med klokka, sentrert i (x,y)
\newcommand{\maske}[3]{%
  \draw[maske] ($(#1,#2)+(60:0.42)$) arc[start angle=60, end angle=-240, radius=0.42];
  \node[text=acc, font=\small] at (#1,#2) {#3};}

% Konvensjoner (testet):
%   V: + ved startpunktet, så en kilde tegnet oppover trenger invert for å få + oppe.
%   I: pila peker i tegneretningen.
%   Loddrett oppover: l = venstre, a = høyre.  Loddrett ned: l_ = venstre, a^ = høyre.
%   Vannrett: l = over, a = under.
% \volt{x}{ytopp}{ybunn}{etikett}: + oppe, - nede, etikett midt mellom, i x
\newcommand{\volt}[4]{%
  \node[font=\small] at (#1,#2) {$+$}; \node[font=\small] at (#1,#3) {$-$};
  \node[font=\small] at (#1,{(#2+#3)/2}) {#4};}
% \volth{y}{xvenstre}{xhoyre}{etikett}: + til venstre, - til høyre
\newcommand{\volth}[4]{%
  \node[font=\small] at (#2,#1) {$+$}; \node[font=\small] at (#3,#1) {$-$};
  \node[font=\small] at ({(#2+#3)/2},#1) {#4};}
\newcommand{\kO}{\,\mathrm{k}\Omega}
\newcommand{\Om}{\,\Omega}

% \thev{ETh}{RTh}: Thévenin-ekvivalent med terminaler a og b, origo nede til venstre
\newcommand{\thev}[2]{%
  \draw (0,0) to[V, invert, l=$E_{Th}$, a=#1] (0,2.6) to[R, l=$R_{Th}$, a=#2] (2.8,2.6)
    to[short, -o] (3.6,2.6) node[above] {$a$};
  \draw (0,0) to[short, -o] (3.6,0) node[below] {$b$};}
% \nort{IN}{RN}: Norton-ekvivalent med terminaler a og b
\newcommand{\nort}[2]{%
  \draw (0,0) to[I, l=$I_N$, a=#1] (0,2.6) -- (2.6,2.6) to[short, -o] (3.6,2.6) node[above] {$a$};
  \draw (2.6,2.6) to[R, l_=$R_N$, a^=#2, *-*] (2.6,0);
  \draw (0,0) to[short, -o] (3.6,0) node[below] {$b$};}
% \last{navn}{verdi}: lastmotstand mellom terminalene a og b (x = 3.6), i aksentfarge
\newcommand{\last}[2]{%
  \draw (3.6,2.6) -- (4.8,2.6) to[R, l_=#1, a^=#2, color=acc] (4.8,0) -- (3.6,0);}

% Transientgrafer (pgfplots): x i tidskonstanter eller sekunder, ingen rutenett
\pgfplotsset{
  trans/.style={width=7.4cm, height=4.6cm, axis lines=left, clip=false, samples=80,
    tick label style={font=\small}, label style={font=\small}, scaled ticks=false,
    /pgf/number format/1000 sep={}, every axis plot/.append style={line width=0.8pt},
    xlabel style={at={(axis description cs:1.03,0)}, anchor=west},
    ylabel style={rotate=-90, at={(axis description cs:0,1.04)}, anchor=south}},
}

% Strømbaner for animasjon (strom.js): f01 ... f12 tegnes i strømretningen og blir
% CSS-klassene flow f01 ... i build.py. Tegnes sist, slik at prikkene ligger over ledningene.
\definecolor{f01}{HTML}{F0F001}
\definecolor{f02}{HTML}{F0F002}
\definecolor{f03}{HTML}{F0F003}
\definecolor{f04}{HTML}{F0F004}
\definecolor{f05}{HTML}{F0F005}
\definecolor{f06}{HTML}{F0F006}
\definecolor{f07}{HTML}{F0F007}
\definecolor{f08}{HTML}{F0F008}
\definecolor{f09}{HTML}{F0F009}
\definecolor{f10}{HTML}{F0F010}
\definecolor{f11}{HTML}{F0F011}
\definecolor{f12}{HTML}{F0F012}
\tikzset{strom/.style={line width=1.6pt, line cap=round, line join=round}}
